\subsection{A finite parameter locus in the degree-ten norm net}
\label{sec:degree-ten-finite-locus}

The first four-pole norm system not excluded by the preceding dimension
tests has a three-dimensional kernel.  The following result makes its
remaining rational-member problem finite; it neither produces a rational
member nor excludes all such members.

Use the Paley word of Appendix~\ref{app:seven-cubic-bank}.  Thus
$\eta^2+\eta+2=0$, $\alpha=(\eta-1)/2$, and
$\gamma=3(\eta+3)/4$.  Its fourteen points in the $(X,Y)$ plane are
\[
 Q_j=(\zeta^j,\zeta^{5j}),\qquad
 T_j=(\alpha\zeta^j,\gamma\zeta^{5j})
 \quad(0\le j<7).
\]
Let $\mathcal N$ consist of polynomials
\[
 F(X,Y)=\sum_{j=0}^{10}A_j(X)Y^j,
 \qquad \deg A_j\le34-3j,
\]
having multiplicity at least four at every $Q_j$ and at least six at
every $T_j$.  Equivalently these are sections of
$10C_0+34f$ on $\mathbb F_3$, with the indicated prescribed points,
where $C_0^2=-3$.

\begin{theorem}\label{thm:degree-ten-finite-rational-locus}
Over characteristic zero, $\dim\mathcal N=3$.
The parameters in $\mathbb P(\mathcal N)$ defining geometrically integral
curves with rational normalization belong to a finite closed set.
An explicit necessary determinantal locus containing them is finite or
empty, and has at most $1\,411\,344$ geometric points.
The same finiteness statement for this necessary locus holds for the
specified reduction modulo $29$.
\end{theorem}

\begin{proof}
There are $220$ coefficient variables and
$7\binom52+7\binom72=217$ prescribed jet conditions.
At the prime with $\eta\mapsto7$ and $\zeta\mapsto16$ modulo $29$,
the exact interpolation matrix has rank $217$.
An exact characteristic-zero construction gives three eigenvectors
$F_0,F_1,F_2$ of characters $1,3,5$ for
$(X,Y)\mapsto(\zeta X,\zeta^5Y)$.
For each character it suffices to solve the $31$ jet conditions at
$(1,1)$ and $(\alpha,\gamma)$ in its $32$ allowed monomials.
Each such matrix has rank $31$.
Normalize the coefficients of $Y^{10},X^2Y^{10},X^4Y^{10}$ in
$F_0,F_1,F_2$ to $25,27,20$, respectively.
These vectors are integral at the chosen prime and reduce to the
archived modular kernel.  The unit $217$-minor proves that they span the
whole kernel, both over characteristic zero and over the local integral
model.  Write $F=aF_0+bF_1+cF_2$.
Every nonzero member retains fiber degree ten: its leading coefficient
is $25a+27bX^2+20cX^4$.
The weighted-infinity restrictions of the three modular basis vectors are
$23z^8+10z$, $15z^9+7z^2$, and $20z^{10}+2z^3$.
Their disjoint character supports and unit coefficients also hold for
the exact lifts.  Thus every nonzero member has actual class
$10C_0+34f$, rather than a smaller class with an artificial infinity fiber.

We first establish a boundary-independent necessary rationality condition.
Let $C$ be an integral member with rational normalization.
Its binary-fiber discriminant $\Delta$ is a section of degree
\[
 18\cdot34-3\cdot10\cdot9=342
\]
on the base.  At any base point, trivialize the ruled surface over the
local DVR.  The binary equation has nonzero reduction, since $C$ has no
fiber component.  After a constant projective change of fiber coordinate,
its leading coefficient is a unit.  The resulting monic order satisfies
the discriminant-index identity
\[
 \operatorname{ord}_x\Delta=r_x+2\delta_x,
\]
where $r_x$ is the total different contribution and $\delta_x$ is the
sum of normalization indices in that fiber.  This argument also covers
fiber-infinite branches and the base point at infinity.

A reduced plane singularity of multiplicity $m$ has
$\delta\ge\binom m2$.  The prescribed points therefore contribute at
least $7\binom42+7\binom62=147$ index units.
Divide $\Delta$ by the fixed degree-$294$ base divisor
\[
 T(X)=\prod_{j=0}^6(X-\zeta^j)^{12}
                    (X-\alpha\zeta^j)^{30}.
\]
Its quotient is a homogeneous binary form $D(X,V)$ of degree $48$.
The divisibility is a polynomial identity throughout the net.
The verified nonzero discriminant gives a dense reduced, separable locus;
the generic leading coefficient is nonzero at the specified base points.
The local order identity and multiplicity--index bound therefore imply
the divisibility there, without requiring ordinary singularities, and
hence identically. Division by the monic fixed polynomial $T$ preserves
the integral coefficient model at the chosen prime.
For a rational degree-ten normalization, Riemann--Hurwitz gives
$\sum_xr_x=18$.  Thus the residual divisor has local orders
$r_x+2k_x$, with $k_x\ge0$ and
\[
 \sum_xk_x=(48-18)/2=15.
\]
At a root of multiplicity $e=r_x+2k_x$, the two homogeneous derivatives
have a common factor of multiplicity at least $e-1\ge k_x$ when
$k_x>0$.  Consequently
\begin{equation}\label{eq:degree-ten-binary-gcd}
 \deg\gcd(D_X,D_V)\ge15.
\end{equation}
This reasoning is also valid in characteristic $29$: the degree-ten
cover is separable and tame, and derivative cancellations can only
increase the common factor.  Since $48$ is a unit, Euler's identity
ensures that a common factor of the derivatives also divides $D$.
A zero discriminant is included in the determinantal condition, but is
not evidence of an integral rational member.

Every discriminant coefficient is homogeneous of degree $18$ in
$(a,b,c)$.  Hence \eqref{eq:degree-ten-binary-gcd} defines a closed locus
$Z$ in the entire projective parameter plane, including $a=0$ and every
base-infinity degree drop.  No boundary or reducible-parameter divisor
is removed.  For example, it is the rank-at-most-$79$ locus of the
$94$-square Sylvester matrix of the degree-$47$ forms $D_X,D_V$.

Here is the finite certificate showing that $Z$ misses the line $c=a$
modulo $29$.  On its affine chart $[a:b:c]=[1:u:1]$, the residual
polynomial has generic degree $48$.  The gcd of its first fifteen low
principal subresultant coefficients with its derivative is, up to a
unit, its degree-$16$ leading coefficient, with factorization
\begin{align*}
 u(u+20)&(u^2+24u+13)
 (u^4+7u^3+5u^2+26u+9)\\
 &\cdot(u^8+7u^7+13u^6+4u^5+17u^4+19u^2+5u+16).
\end{align*}
Away from these factors, the subresultants exclude
\eqref{eq:degree-ten-binary-gcd}.  At their roots, computation over the
respective residue fields of degrees $1,1,2,4,8$ gives residual degree
$47$ and derivative-gcd degrees $0,2,0,0,0$.
The root at base infinity then has multiplicity one and contributes
nothing to the homogeneous derivative gcd.
At the remaining projective parameter $[0:1:0]$, the residual degree is
$45$, its finite derivative gcd has degree two, and base infinity
contributes two, giving total four.  Thus the entire projective line
misses $Z$, including all extension-field points and all leading-degree
drops.  None of these tested discriminants is zero.

Every positive-dimensional closed subset of $\mathbb P^2$ meets every
projective line.  Therefore $Z$ is finite or empty modulo $29$.
The integral kernel model gives the same projective determinantal
locus and the line $c=a$ over the local DVR.  Their intersection is
proper over that DVR.  A geometric generic point in the intersection
would specialize, after finite base extension, to a geometric special
point.  Since the latter does not exist, the characteristic-zero locus
also misses the line and is finite or empty.

For an explicit coarse bound, consider the structured map
\[
 \operatorname{Sym}^{32}(k^2)\oplus\operatorname{Sym}^{32}(k^2)
 \longrightarrow\operatorname{Sym}^{79}(k^2),
 \qquad (A,B)\longmapsto A D_X+B D_V.
\]
Its $80$-by-$66$ matrix has entries of parameter degree $18$.
For two degree-$47$ forms with gcd degree $d$, the primitive syzygy has
degree $47-d$; hence the matrix has a kernel precisely when $d\ge15$.
This also includes a zero form.  Its maximal minors, of degree
$66\cdot18=1188$, thus define a determinantal scheme with support $Z$.
Because this support is finite, two general linear combinations of the
minors have no common curve.  Their complete intersection contains the
scheme and has length $1188^2=1\,411\,344$, proving the claimed bound.
\end{proof}

\paragraph{Exact certificates and reproduction.}
The accompanying directory
\path{research/prime_line_strengthening/degree_ten_four_pole_route}
contains the coefficient matrix and unit minor in \texttt{gate.json},
the three exact $\mathbb Q(\eta)$ eigenvectors in
\texttt{exact\_kernel.json}, and independent replays of both calculations.
In its \texttt{residual\_discriminant} subdirectory,
\texttt{reconstruction.json} stores every coefficient of
$D(X,1;1,s,t)$.  Total parameter degree $18$ permits exact triangular
Newton interpolation from the $190$ pairs
$s,t\ge0$, $s+t\le18$ modulo $29$; all relevant factorials are units.
Each evaluation computes the degree-ten discriminant and divides by
$T$ exactly.  Independent dense interpolation and additional resultant
evaluations verify the reconstruction.  Homogenizing its coefficients
to parameter degree $18$ recovers the entire parameter plane.
The complete line subresultants and extension-field checks are in
\texttt{line-t1.json}; the projective endpoint is independently replayed
in \texttt{endpoint\_independent.json}.

The theorem concerns a necessary locus in this particular Paley norm net.
Its isolated points have not been enumerated or classified here, and
the theorem supplies no new nearby polynomial or augmentation.
